\chapter{Ejercicios de Decibilidad}

\setcounter{ejercicio}{7}

% Ejercicio 8

\begin{ejercicio}
    Demostrar que el problema de la parada, determinar el conjunto de parejas $(M, w)$ tales
    que la MT $M$ para cuando tiene a $w$ como entrada, es semidecidible pero no decidible.
\end{ejercicio}

% Ejercicio 9

\begin{ejercicio}
    (MANUAL) Demostrar que determinar si una MT acepta al menos un palíndromo no es
    decidible pero sí semidecidible.
\end{ejercicio}

% Ejercicio 10

\begin{ejercicio}
    Supongamos que tenemos MT con dos tipos de estados: \textit{campana} y \textit{silbato}. Determinar que saber si una MT entrará alguna vez en un estado \textit{silbato} es indecidible. ¿Es semidecidible?
\end{ejercicio}

% Ejercicio 11

\begin{ejercicio}
    (MANUAL) Demostrar que es indecidible (no recursivo) saber si una MT termina escribiendo un 1 cuando comienza con una cinta completamente en blanco.
\end{ejercicio}

% Ejercicio 12

\begin{ejercicio}
    Determinar si los siguientes problemas son decidibles, semidecidibles o no semidecidibles:
    \begin{itemize}
        \item[a)] Determinar si el lenguaje de una MT contiene, al menos, dos palabras distintas.
        \item[b)] Determinar si el lenguaje reconocido por una MT es finito.
        \item[c)] (MANUAL) Determinar si el lenguaje reconocido por una MT es infinito.
        \item[d)] Determinar si el lenguaje de una MT es independiente del contexto.
    \end{itemize}
\end{ejercicio}

% Ejercicio 13

\begin{ejercicio}
    Sea $L$ el lenguaje formado por pares de códigos de MT más un entero $(M_1 ,M_2 , k)$ tales que $L(M_1) \cap L(M_2)$ contiene, al menos, $k$ palabras. Demostrar que $L$ es semidecidible pero no decidible.
\end{ejercicio}

% Ejercicio 14

\begin{ejercicio}
    Demostrar que los siguientes lenguajes son recursivos:
    \begin{itemize}
        \item[a)] El conjunto de las MT $M$ tales que al comenzar con la cinta en blanco, en algún momento escribirán un símbolo no blanco en la cinta.
        \item[b)] El conjunto de las MT que nunca se mueven a la izquierda.
        \item[c)] El conjunto de los pares $(M, w)$ tales que $M$, al actuar sobre la entrada $w$, nunca lee una casilla de la cinta más de una vez.
    \end{itemize}
\end{ejercicio}

% Ejercicio 15

\begin{ejercicio}
    Indicar si los siguientes problemas son decidibles, semidecidibles o no semidecidibles:
    \begin{itemize}
        \item[a)] (MANUAL) Determinar si una MT se para para cualquier entrada.
        \item[b)] Determinar si una MT no se para para ninguna entrada.
        \item[c)] Determinar si una MT se para, al menos, para una entrada.
        \item[d)] Determinar si una MT no se para, al menos, para una entrada.
    \end{itemize}
\end{ejercicio}

% Ejercicio 16

\begin{ejercicio}
    Supongamos que tenemos $4$ lenguajes $L_1$, $L_2$, $L_3$, $L_4$ de tal manera que existe una reducción de $L_1$ a $L_2$, de $L_2$ a $L_3$ y de $L_4$ a $L_3$. Para cada una de las siguientes afirmaciones,
    indicar si para todos los casos son CIERTAS, para todos los casos son FALSAS o si son POSIBLES (en algunos casos son ciertas y, en otros, son falsas).
    \begin{itemize}
        \item[a)] $L_1$ es recursivamente enumerable pero no recursivo, y $L_3$ es recursivo. 
        \item[b)] $L_1$ no es recursivo y $L_4$ no es recursivamente enumerable.
        \item[c)] El complementario de $L_1$ no es recursivamente enumerable, pero el complementario de $L_2$ es recursivamente enumerable.
        \item[d)] El complementario de $L_2$ no es recursivo, pero el complementario de $L_3$ es recursivo.
        \item[e)] Si $L_1$ es recursivo, entonces el complementario de $L_2$ es recursivo.
        \item[f)] Si $L_3$ es recursivo, entonces el complementario de $L_4$ es recursivo.
        \item[g)] Si $L_3$ es recursivamente enumerable, entonces la unión de $L_2$ y $L_4$ es recursivamente enumerable.
        \item[h)] Si $L_3$ es recursivamente enumerable, entonces la intesección de $L_2$ y $L_4$ es recursivamente enumerable.
    \end{itemize}
\end{ejercicio}

% Ejercicio 17

\begin{ejercicio}
    Sea el problema de determinar si una máquina de Turing acepta a lo más 100 palabras. Determinar si es decidible, semidecidible o no semidecidible. Justificar la respuesta.
\end{ejercicio}

% Ejercicio 18

\begin{ejercicio}
    Determinar cuáles de los siguientes problemas son decidibles, semidecidibles o no semidecidibles:
    \begin{itemize}
        \item[a)] Saber si una MT para una entrada 0011 no va a usar más de 10 casillas de la cinta.
        \item[b)] Dadas dos MT saber si aceptan el mismo lenguaje
    \end{itemize}
\end{ejercicio}

% Ejercicio 19

\begin{ejercicio}
    Determinar cuáles de los siguientes problemas son decidibles, semidecidibles o no semidecidibles:
    \begin{itemize}
        \item[a)] Dada una MT $M$ y una palabra $u$, saber si la MT acepta la palabra $u$ en un número de pasos menor o igual a $|u|$.
        \item[b)] Dada una MT determinar si no acepta la palabra vacía.
    \end{itemize}
\end{ejercicio}

% Ejercicio 20

\begin{ejercicio}
    Determinar cuáles de los siguientes problemas son decidibles, semidecidibles o no semidecidibles 
    (se supone que las MTs tienen a \{0, 1\} como alfabeto de entrada):
    \begin{itemize}
        \item[a)] Dadas dos máquinas de Turing, $M_1$ y $M_2$, determinar si existe, al menos, una palabra
        aceptada simultáneamente por ambas máquinas.
        \item[b)] (MANUAL) Dada una MT, determinar si para toda palabra de entrada no realiza más de 5 movimientos.
        \item[c)] Determinar si una MT no es determinista.
        \item[d)] (MANUAL) Dada una MT, determinar si ninguna de las palabras que acepta es un palíndromo.
    \end{itemize}
    Justifica las respuestas.
\end{ejercicio}

% Ejercicio 21

\begin{ejercicio}
    Se considera el siguiente problema: dadas dos máquinas de Turing, determinar si aceptan el mismo lenguaje. 
    Razonar si este problema es decidible, semidecidible o no semidecidible.
\end{ejercicio}

% Ejercicio 22

\begin{ejercicio}
    Determinar, justificando las respuestas, cuáles de los siguientes problemas son decidibles, semidecidibles o no semidecidibles:
    \begin{itemize}
        \item[a)] Dada una máquina de Turing y una palabra de entrada determinar si visita menos de 10 casillas distintas de la cinta de lectura. % Decidible
        \item[b)] Dadas dos Máquinas de Turing, determinar si el conjunto de las palabras aceptadas a la vez por ambas máquinas de Turing es finito. % No semidecidible
    \end{itemize}
\end{ejercicio}

% Ejercicio 23

\begin{ejercicio}
    Indica cuáles de los siguientes problemas son decidibles, semidecidibles o no semidecidibles:
    \begin{itemize}
        \item[a)] Determinar si el lenguaje aceptado por una MT es regular. % No decidible
        \item[b)] Dadas dos MTs determinar si hay una palabra que es aceptada por las dos MT y además es un palíndromo. % Semidecidible y no decidible
    \end{itemize}
\end{ejercicio}

% Ejercicio 24

\begin{ejercicio}
    Indica cuáles de los siguientes problemas son decidibles, semidecidibles o no semidecidibles:
    \begin{itemize}
        \item[a)] Dada una MT determinar si todos sus estados son accesibles (se puede llegar a ellos para un cálculo para alguna entrada).
        \item[b)] Dada una MT determinar si su número de estados es menor o igual a 5.
    \end{itemize}
\end{ejercicio}

% Ejercicio 25

\begin{ejercicio}
    Demuestra que el Problema de la Correspondencia de Post con un alfabeto A que tiene un sólo elemento es decidible.
\end{ejercicio}

\newpage
\section{Soluciones}
\setcounter{ejercicio}{7}

% Ejercicio 8

\begin{ejercicio}
    Demostrar que el problema de la parada, determinar el conjunto de parejas $(M, w)$ tales
    que la MT $M$ para cuando tiene a $w$ como entrada, es semidecidible pero no decidible. \\

    Definimos $\text{PARADA}(M,w) = \{\cod{M}{w} : M \text{ para con entrada } w\}$. Para ver que es 
    semidecidible, la idea es construir un simulador de MT $M_p$, que lea la codificación
    $\cod{M}{w}$, simule la MT $M$ sobre la entrada $w$ y si en algún momento de la simulación $M$ para, $M_p$
    acepta, de manera que $M_p$ acepte exactamente las entradas $\cod{M}{w}$ tal que $M$ para con $w$.

    \begin{algorithm}
        \caption{Máquina de Turing $M_p$}
        \label{algorithm:parada-semidecidible}
        \begin{algorithmic}
            \State \textbf{Entrada:} Una palabra $v \in \{0,1\}^*$

            \State Comprobar si $v$ es de la forma $\langle M,w\rangle$
            \If{$v$ no es una codificación válida}
                \State Rechazar
            \EndIf

            \State Simular paso a paso la ejecución de $M$ sobre la entrada $w$

            \If{en algún momento la simulación de $M$ para}
                \State Aceptar
            \EndIf
        \end{algorithmic}
    \end{algorithm}

    Esta MT $M_p$ acepta exactamente los casos positivos de las palabras del lenguaje $L_p$ asociado:
    \begin{itemize}
        \item Si $\cod{M}{w}\in \text{PARADA}$, la simulación de $M$ sobre $w$ termina en un número finito de pasos, y 
        $M_p$ aceptará.
        \item Si $\cod{M}{w}\notin \text{PARADA}$, entonces $M$ cicla indefinidamente sobre $w$, y la simulación hace continúa
        también indefinidamente, por tanto $M_p$ no aceptará.
    \end{itemize}

    Por lo tanto, PARADA es semidecidible. \\

    Para ver que no es decidible, la idea será hacer una reducción de universal a parada, que sabemos que es semidecidible, 
    pero no decidible por teoría. Consideramos los problemas:
    $$\begin{cases}
        P1 = &\text{UNIVERSAL : dada una MT $M$ y una palabra $w$, determinar si $M$ acepta $w$} \\
        P2 = &\text{PARADA : dada una MT $N$ y una palabra $u$, determinar si $N$ para con $u$}
    \end{cases}$$

    

    Construimos un algoritmo $F$ que transforma cada instancia $(M,w)$ de UNIVERSAL en una instancia $(N,w)$
    de PARADA, donde $N = F(M)$ es la MT siguiente:

    \begin{algorithm}
        \caption{Máquina de Turing $N = F(M)$}
        \label{algorithm:reduccion-universal-parada}
        \begin{algorithmic}
            \State \textbf{Entrada:} Una palabra $u \in \{0,1\}^*$

            \State 1:\ Simular la ejecución de $M$ sobre la entrada $u$

            \State 2:\ Si $M$ acepta $u$, entonces $N$ para

            \State 3:\ En otro caso, $N$ no para (cicla).
        \end{algorithmic}
    \end{algorithm}

    Esta construcción es algorítmica, ya que a partir de $M$ puede construirse una MT $N$ que simule a $M$ sobre su 
    misma entrada $u$ y que solo pare cuando la simulación llegue a aceptación. \\

    Veamos la equivalencia:

    \begin{itemize}
        \item Si $(M,w)$ es un caso positivo de UNIVERSAL, es decir, si $M$ acepta $w$, entonces al ejecutar $N$ sobre
        $w$ la simulación llega a aceptación y $N$ para. Por tanto, $(N,w) \in \text{PARADA}$.
        \item Si $(M,w)$ es un caso negativo de UNIVERSAL, es decir, si $M$ no acepta $w$, entonces la máquina $N$,
        sobre $w$, no para. Por tanto, $(N, w) \notin \text{PARADA}$. Así 
        $$(M,w) \in \text{UNIVERSAL} \iff (N,w) \in \text{PARADA}$$
    \end{itemize}

    Por tanto, tenemos que UNIVERSAL$(M,w)$ $\propto$ PARADA$(N,w)$. Si PARADA fuera decidible, tendríamos que UNIVERSAL también
    lo sería porque UNIVERSAL se reduce a PARADA. Sin embargo, sabemos que esto no es cierto, por lo que PARADA
    no es decidible.
\end{ejercicio}



% Ejercicio 9

\begin{ejercicio}\label{ejercicio:ej9}
    (MANUAL) Demostrar que determinar si una MT acepta al menos un palíndromo no es
    decidible pero sí semidecidible. \\

    Dada una MT $M$, definimos $\text{PAL}(M) = 
    \{\langle M \rangle : \exists w \in \{0,1\}^{*} \text{ palíndromo tal que } M \text{ acepta } w\}$. \\

    Para la no decidibilidad, vemos que la propiedad ``aceptar al menos un palíndromo'' depende solo del lenguaje aceptado
    por la máquina, $L(M)$, y no de la máquina concreta, es decir, es una propiedad de los lenguajes r.e.. Además, 
    es una propiedad no trivial, ya que el lenguaje $\{11\}$ la cumple, pero $\{10\}$ no. Así, por el Teorema de Rice, 
    PAL no es decidible. \\

    Para la semidecidibilidad, consideramos una MTND $M_{PAL}$ que funciona como sigue:

    \begin{algorithm}
        \caption{MTND $M_{PAL}$}
        \label{algorithm:pal-semidecidible}
        \begin{algorithmic}
            \State \textbf{Entrada:} Una codificación $\langle M \rangle$ de una MT $M$

            \State Selecciona de forma no determinista una palabra $w \in \{0,1\}^{*}$.

            \State Comprueba si $w$ es un palíndromo. 
            
            \State Si $w$ es palíndromo, simula $M$ con entrada $w$

            \State Si $M$ acepta $w$, entonces $M_{PAL}$ acepta.
        \end{algorithmic}
    \end{algorithm}

    Tenemos entonces que, si $M$ acepta al menos un palíndromo $w_0$, existe una rama de cómputo de $M_{PAL}$ que 
    precisamente toma $w_0$, comprueba que es palíndromo y acepta. Sin embargo, si $M$ no acepta ningún palíndromo, 
    entonces ninguna rama puede aceptar. Por lo tanto, $M_{PAL}$ acepta exactamente las codificaciones $\langle M \rangle$
    tales que $M$ acepta al menos un palíndromo, de donde PAL$(M)$ es semidecidible.

\end{ejercicio}



% Ejercicio 10

\begin{ejercicio}
    Supongamos que tenemos MT con dos tipos de estados: \textit{campana} y \textit{silbato}. 
    Determinar que saber si una MT entrará alguna vez en un estado \textit{silbato} es indecidible. ¿Es semidecidible? \\

    Definimos $$\text{SILBATO}(M) = \{\langle M \rangle : \exists w \in \{0,1\}^{*} : M \text{ entra alguna vez en un estado silbato con entrada } w\}$$

    Para ver que no es decidible, dada una MT $M$, como la propiedad ``entrar en un estado silbato'' 
    no es una propiedad del lenguaje $L(M)$, tampoco es una propiedad de los lenguajes r.e., y no puede aplicarse el 
    Teorema de Rice. \\

    Haremos entonces una reducción desde un problema no decidible de tal manera que sea sencilla. Por ejemplo, PARADA.
    Consideramos entonces el problema $$\text{PARADA}(M,w) = \{\langle M, w \rangle : M \text{ para con entrada } w\}$$

    Vamos a construir un proceso algorítmico $F$ para pasar de una instancia de PARADA a una instancia de 
    $SIL$. Sea $(M, w)$ una instancia de PARADA, construimos las siguiente máquina de Turing $F(M, w)$ 
    (que es una instancia de SIL)

    \begin{algorithm}
        \caption{MT $F(M,w)$}
        \label{algorithm:sil-indecidible}
        \begin{algorithmic}
            \State \textbf{Entrada:} $v \in \{0,1\}^{*}$

            \State Borramos $v$ de la cinta.
            \State Copiamos $\langle M \rangle 111w$ en la cinta.

            \State Simulamos $M$ con entrada $w$. 

            \State \hspace{1cm} \textbf{Si} $M$ para \textbf{entonces}

            \State \hspace{2cm} $F(M,w)$ entra en un estado \textit{silbato}.

            \State \hspace{1cm} \textbf{Si} $M$ no para \textbf{entonces}

            \State \hspace{2cm} $F(M,w)$ cicla y nunca entra en un estado \textit{silbato}.
        \end{algorithmic}
    \end{algorithm}

    Vemos que la máquina $F(M,w)$ se construye de tal manera que todos los estados usados para las tareas previas a que 
    $M$ termine sean de tipo \textit{campana}. Si $M$ termina, entonces añadimos un nuevo estado especial $q_s$, de tipo
    \textit{silbato}, al que se llega únicamente cuando la simulación termina. Por lo tanto
    $$\langle M,w \rangle \in \text{PARADA} \iff \langle F(M,w) \rangle \in \text{SILBATO}$$

    Puesto que:
    \begin{enumerate}
        \item Si $M$ para con $w$ (caso positivo de PARADA), entonces $F(M,w)$ termina la simulación y entra en $q_s$ 
        (caso positivo de SILBATO).
        \item Si $M$ no para con $w$ (caso negativo de PARADA), entonces $F(M,w)$ ciclará y nunca 
        entrará en estado silbato (caso negativo de SILBATO).
    \end{enumerate}

    Por lo tanto, PARADA $\propto$ SILBATO. Si SILBATO fuera decidible, también lo sería PARADA, cosa que sería una contradicción.
    Así, SILBATO no es decidible. \\

    La semidecidibilidad puede comprobarse con una MTND que seleccione de forma no determinista una entrada $u$ y simule $M$
    con dicha entrada y acepte en caso de que durante la simulación $M$ entre en un estado \textit{silbato}. Más en detalle:

    \begin{algorithm}
        \caption{MTND $M_{SILBATO}$}
        \label{algorithm:sil-semidecidible}
        \begin{algorithmic}
            \State \textbf{Entrada:} Una codificación $\langle M \rangle$ de una MT $M$
                \State Selecciona de forma no determinista una palabra $u \in \{0,1\}^{*}$.

                \State Simulamos $M$ con entrada $u$.

                \State Si en algún momento de la simulación $M$ entra en un estado \textit{silbato}, entonces 
                $M_{SILBATO}$ acepta.
        \end{algorithmic}
    \end{algorithm}
\end{ejercicio}



% Ejercicio 11

\begin{ejercicio}
    (MANUAL) Demostrar que es indecidible (no recursivo) saber si una MT termina escribiendo un 1 cuando comienza con 
    una cinta completamente en blanco. \\

    Hecho en el manual. Para la no decidiblidad hay que hacer una reducción desde UNIVERSAL al problema, y para la 
    semidecidibilidad basta ejecutar un bucle $i =1, \ldots, \infty$ tal que se simule la MT con entrada $\varepsilon$ en
    $i$ pasos, y si en alguno de esos $i$ pasos escribe un $1$ acepta la entrada.
\end{ejercicio}



% Ejercicio 12

\begin{ejercicio}
    Determinar si los siguientes problemas son decidibles, semidecidibles o no semidecidibles:
    \begin{itemize}
        \item[a)] Determinar si el lenguaje de una MT contiene, al menos, dos palabras distintas.\\
        
        Definimos $\text{DOSPal}(M) = \{\langle M \rangle : \exists u,v \in L(M) \text{ con } u \neq v\}$. Vemos que 
        no es decidible, ya que la propiedad ``$L(M)$ contiene al menos dos palabras distintas'' depende solo
        del lenguaje $L(M)$, no de la MT $M$ concreta, luego es una propiedad de los lenguajes r.e.. Además, es no trivial,
        ya que $\{0\}$ no cumple la propiedad, pero $\{0,1\}$ sí la cumple. Por el Teorema de Rice, DOSPal$(M)$ es no
        decidible. Veamos que es semidecidible. Para ello, consideramos una MTND $M_{DOSPal}$ que funciona como sigue:

        \begin{algorithm}
            \caption{MTND $M_{DOSPAL}$}
            \label{algorithm:dospal-semidecidible}
            \begin{algorithmic}
                \State \textbf{Entrada:} Una codificación $\langle M \rangle$ de una MT $M$

                \State Selecciona de forma no determinista dos palabras $u,v \in \{0,1\}^{*}$.

                \State Si $u = v$, entonces $M_{DOSPal}$ rechaza.

                \State Simulamos $M$ con entrada $u$ y con entrada $v$. % no hace falta que sean paralelo

                \State Si ambas simulaciones aceptan, entonces $M_{DOSPal}$ acepta.
            \end{algorithmic}
        \end{algorithm}

        Tenemos entonces que, si $M$ acepta al menos dos palabras distintas $u_0$, $v_0$, existe una rama de cómputo de 
        $M_{DOSPal}$ que precisamente toma $u_0$ y $v_0$, simula $M(u_0)$ y $M(v_0)$, y como ambas aceptan, esa rama acepta. 
        Sin embargo, si $M$ no acepta dos palabras distintas, entonces para cualquier par distinto $(u,v)$, al menos una de 
        las dos simulaciones no aceptará y consecuentemente esa rama no aceptará. Por lo tanto, $M_{DOSPal}$ acepta 
        exactamente las codificaciones $\langle M \rangle$ tales que $M$ acepta al menos dos palabras distintas, 
        de donde DOSPal$(M)$ es semidecidible.

        

        \item[b)] Determinar si el lenguaje reconocido por una MT es finito. \\
        
        Definimos $\text{FINITO}(M) = \{\langle M \rangle : L(M) \text{ es finito}\}$. Vemos que 
        no es decidible, ya que la propiedad ``$L(M)$ es finito'' depende solo del lenguaje $L(M)$, 
        no de la MT $M$ concreta, luego es una propiedad de los lenguajes r.e.. Además, es no trivial,
        ya que $\empty$ cumple la propiedad (una MT sin estado final), pero $\{0,1\}^{*}$ no la cumple
        (una MT con el estado inicial también final que acepte todas las palabras). Por el Teorema de Rice, FINITO$(M)$ es no
        decidible. \\

        Para lo no semidecidibilidad se puede hacer una reducción desde $$\text{DIAGONAL}(M) = \{\langle M \rangle : 
        M \text{ no acepta } \langle M \rangle \}$$ 
        (que sabemos por teoría que no es semidecidible), igual que en el apartado c), pero permutando los casos en 
        que se rechaza por lo que se acepta. El proceso algorítmico $F$ pasará de una instancia de DIAGONAL a una 
        instancia de FINITO. Sea entonces $M$ una instancia de DIAGONAL, y construimos la siguiente máquina de Turing $F(M)$:

        \begin{algorithm}
            \caption{MT $F(M)$}
            \label{algorithm:finito-no-semidecidible}
            \begin{algorithmic}
                \State \textbf{Entrada:} Una palabra $v \in \{0,1\}^{*}$.

                \State \hspace{1cm} \textbf{Si} $v$ no es de la forma $0^n$ con $n > 0$ \textbf{entonces} 

                \State \hspace{2cm} $F(M)$ rechaza.

                \State \hspace{1cm} \textbf{Si} $v = 0^n$ para cierto $n > 0$ \textbf{entonces} 

                \State \hspace{2cm} Simulamos $M$ con entrada $\langle M \rangle$ durante $n$ pasos.

                \State \hspace{2cm} \textbf{Si} $M$ acepta $\langle M \rangle$ \textbf{entonces}

                \State \hspace{3cm} $F(M)$ acepta.

                \State \hspace{2cm} \textbf{Si no} \textbf{entonces}

                \State \hspace{3cm} $F(M)$ rechaza.
            \end{algorithmic}
        \end{algorithm}

        \item[c)] (MANUAL) Determinar si el lenguaje reconocido por una MT es infinito. \\
        
        Hecho en el manual.

        \item[d)] Determinar si el lenguaje de una MT es independiente del contexto. \\
        
        Definimos $\text{INDCONT}(M) = \{\langle M \rangle : L(M) \text{ es independiente del contexto}\}$. Al menos es no 
        decidible por el Teorema de Rice, ya que, para una MT $M$, la propiedad ``ser $L(M)$ 
        independiente del contexto'' es una propiedad que depende solo de $L(M)$, no de $M$, luego es una propiedad
        no trivial de los lenguajes r.e. (no trivial porque sabemos que existe una MT $M_1$ tal que $L(M_1)$ es independiente 
        del contexto y $M_2$ tal que $L(M_2)$ no es independiente del contexto (por Modelos de Computación)). De hecho, podemos ver 
        que es no semidecidible por medio de una reducción desde
        $$\text{DIAGONAL}(M) = \{\langle M \rangle : M \text{ no acepta } \langle M \rangle \}$$

        Sea $M''$ una MT fija tal que
        $$
        L(M'')=\{0^n1^n0^n : n\geq 0\},
        $$
        que no es un lenguaje independiente del contexto (puede comprobarse con el Lema de bombeo). Dada una MT $M$, construimos una MT $M'$ con el siguiente comportamiento:

        \begin{algorithm}
            \caption{MT $M'$}
            \label{algorithm:reduccion-diagonal-indcont}
            \begin{algorithmic}
                \State \textbf{Entrada:} Una palabra $x \in \{0,1\}^{*}$

                \State Simulamos $M$ con entrada $\langle M\rangle$.

                \If{$M$ acepta $\langle M\rangle$}
                    \State Simulamos $M''$ con entrada $x$.
                    \If{$M''$ acepta $x$}
                        \State $M'$ acepta.
                    \Else
                        \State $M'$ rechaza.
                    \EndIf
                \EndIf
            \end{algorithmic}
        \end{algorithm}

        Tenemos entonces que:

        \begin{itemize}
            \item Si $\langle M\rangle \in \text{DIAGONAL}$, entonces $M$ no acepta
                $\langle M\rangle$. Por tanto, al simular $M$ con entrada $\langle M\rangle$,
                pueden pasar dos cosas: o bien $M$ rechaza, o bien $M$ cicla. En ambos casos,
                $M'$ no acepta ninguna palabra. Luego $$L(M')=\emptyset.$$
                Como $\emptyset$ es independiente del contexto, se tiene que
                $$\langle M'\rangle \in \text{INDCONT}.$$
            \item Si $\langle M\rangle \notin \text{DIAGONAL}$, entonces $M$ acepta
                $\langle M\rangle$. Por construcción, en ese caso $M'$ pasa a comportarse como
                $M''$ sobre su entrada $x$. Por tanto, $$L(M')=L(M'')=\{0^n1^n0^n : n\geq 0\}.$$
                Como este lenguaje no es independiente del contexto, se tiene que
                $$\langle M'\rangle \notin \text{INDCONT}.$$
        \end{itemize}

        Así, $$\langle M\rangle \in \text{DIAGONAL} \Longleftrightarrow \langle M'\rangle \in \text{INDCONT}.$$
        Por tanto, $$ \text{DIAGONAL} \propto \text{INDCONT}.$$
        Como $\text{DIAGONAL}$ es no semidecidible, $\text{INDCONT}$ es no semidecidible.
    \end{itemize}
\end{ejercicio}



% Ejercicio 13

\begin{ejercicio}
    Sea $L$ el lenguaje formado por pares de códigos de MT más un entero $(M_1 ,M_2 , k)$ tales que 
    $L(M_1) \cap L(M_2)$ contiene, al menos, $k$ palabras. Demostrar que $L$ es semidecidible pero no decidible. \\

    Vemos que $$L = \{(M_1, M_2, k) : |L(M_1) \cap L(M_2)| \geq k\}$$

    Para ver que es semidecidible, construimos la siguiente $MTND$:

    \begin{algorithm}
        \caption{MTND $M_L$}
        \label{algorithm:L-semidecidible}
        \begin{algorithmic}
            \State \textbf{Entrada:} Una codificación $\langle M_1,M_2,k\rangle$, donde $M_1$ y $M_2$ son MT y $k\in\mathbb{Z}$.

            \If{$k=0$}
                \State $M_L$ acepta.
            \EndIf

            \State Selecciona de forma no determinista $k$ palabras $u_1,\dots,u_k \in \{0,1\}^{*}$.

            \If{$\exists i,j \in \{1,\dots,k\} : i\neq j \quad \land \quad u_i=u_j$}
                \State $M_L$ rechaza.
            \EndIf

            \For{$i=1,\dots,k$}
                \State Simulamos $M_1$ con entrada $u_i$.
                \State Simulamos $M_2$ con entrada $u_i$.

                \If{alguna de las dos simulaciones rechaza}
                    \State $M_L$ rechaza.
                \EndIf
            \EndFor

            \If{todas las simulaciones aceptan}
                \State $M_L$ acepta.
            \EndIf
        \end{algorithmic}
    \end{algorithm}

    Vemos que funciona, ya que $(M_1, M_2, k) \in L$ si y solo si existen $k$ palabras distintas $u_1, \ldots, u_k$ tales
    que $u_i \in L(M_1) \cap L(M_2) \quad \forall i = 1, \ldots, k$. Por lo tanto, $M_L$ semidecide $L$, luego 
    $L$ es semidecidible. \\

    Para ver que $L$ es no decidible, hacemos una reducción desde 
    $$\text{C-VACIO}(M) = \{\langle M \rangle : L(M) \neq \emptyset\}$$
    que sabemos que no es decidible por teoría. \\

    Dada una MT $M$, construimos la instancia $(M, U, 1)$, con $U$ es una $MT$ que acepta todas las palabras, es decir,
    $L(U) = \{0,1\}^{*}$ (por ejemplo, $U$ tiene el estado inicial también final). Entonces 
    $$(M,U,1) \in L \iff |L(M) \cap L(U)| \geq 1 \iff L(M) \neq \emptyset$$
    donde en la última equivalencia se ha usado que $L(M) \cap L(U) = L(M)$, y $|L(M)| \geq 1 \iff L(M) \neq \emptyset$. 
    Tenemos entonces que $$M \in \text{C-VACIO} \iff (M,U,1) \in L$$
    Así, C-VACIO $\propto L$, y si $L$ fuera decidible también lo sería C-VACIO, cosa que es falsa, luego $L$ no es decidible.
\end{ejercicio}



% Ejercicio 14

\begin{ejercicio}
    Demostrar que los siguientes lenguajes son recursivos:
    \begin{itemize}
        \item[a)] El conjunto de las MT $M$ tales que al comenzar con la cinta en blanco, en algún momento 
        escribirán un símbolo no blanco en la cinta. \\

        Sea $$\hspace{-1cm}L_a = \{\langle M \rangle : M \text{ empezando con la cinta en blanco escribe alguna vez un símbolo distinto 
        de \#}\}$$

        Teniendo en cuenta que, si la cinta comienza entera en blanco (entonces $M$ se ejecuta sobre la palabra vacía), 
        seguirá estándolo mientras no se escriba ningún símbolo no blanco. Por tanto, 
        en cada paso $M$ lee siempre \#. El comportamiento de $M$ se puede determinar
        entonces por transiciones de la forma $\delta(q,\#)$. Como una MT tiene un conjunto finito de estados 
        y transiciones de la forma $\delta(q,b) = (p, c, M)$, con $M \in \{I,S,D\}$, puede decidirse algorítmicamente
        en una cantidad finita de pasos, y por tanto en tiempo finito.

        \begin{algorithm}
            \caption{MT $M_{L_a}$}
            \begin{algorithmic}
            \State \textbf{Entrada:} Una codificación $\langle M\rangle$ de una MT $M$
            \State $q \gets q_0$
            \State $V \gets \emptyset$
            \While{$q \notin V$}
                \State Añadir $q$ a $V$
                \If{$\delta(q,\#)$ no está definida}
                    \State Rechazar
                \EndIf
                \State Sea $\delta(q,\#)=(p,c,m)$
                \If{$c\neq \#$}
                    \State Aceptar
                \EndIf
                \State $q \gets p$
            \EndWhile
            \State Rechazar
            \end{algorithmic}
        \end{algorithm}

        Si aparece una transición que escribe un símbolo $c \neq \#$, $M$ escribe un símbolo no blanco, luego aceptamos.
        En otro caso, la máquina nunca escribirá un símbolo no blanco. Como los estados son finitos, el algoritmo
        siempre termina, y decide correctamente.

        

        \item[b)] El conjunto de las MT que nunca se mueven a la izquierda. \\
        
        Sea 
        $$L_b = \{\langle M \rangle : M \text{ nunca se mueve a la izquierda}\}$$

        Para ver que es decidible, basta con buscar si hay alguna transición con movimiento $I$ que pueda ejecutarse.
        Como antes de la primera vez la máquina se mueve a la izquierda, todos los movimientos son estáticos o a la derecha,
        la máquina nunca vuelve a leer una casilla anterior. Por tanto, antes de ese posible primer movimiento a la izquierda,
        lo único que puede leer es, o bien, símbolos de la entrada, o bien, si los lee todos, blancos \#. \\

        Así, podemos construir un grafo finito de configuraciones resumidas, con tripletas $(q,t, b)$, con $q \in Q$ el estado 
        actual, $t \in \{\text{entrada, blancos}\}$ que indica si aún estamos leyendo parte de la entrada, o ya hemos
        terminado con la entrada y hemos llegado a blanco, y $b \in B$ que indica el símbolo leído en la casilla actual. \\

        \begin{algorithm}
            \caption{Decisor para $L_b$}
            \begin{algorithmic}
                \State \textbf{Entrada:} Una codificación $\langle M\rangle$ de una MT $M$.

                \State Construimos un grafo finito cuyos nodos son ternas $(q,t,b)$.

                \State Los nodos iniciales son $(q_0,\text{entrada},a)$, para cada $a\in A$, y $(q_0,\text{blancos},\#)$.

                \For{cada nodo $(q,t,b)$}
                    \If{$\delta(q,b)=(p,c,I)$}
                        \State Marcamos $(q,t,b)$ con (*).
                    \ElsIf{$\delta(q,b)=(p,c,S)$}
                        \State Añadimos una arista dirigida de $(q,t,b)$ a $(p,t,c)$.
                    \ElsIf{$\delta(q,b)=(p,c,D)$}
                        \If{$t=\text{entrada}$}
                            \State Añadimos aristas dirigidas a $(p,\text{entrada},a)$, para cada $a\in A$.
                        \EndIf
                        
                        \State Añadimos una arista dirigida a $(p,\text{blancos},\#)$.
                    \EndIf
                \EndFor

                \State Si algún nodo marcado con (*) es alcanzable desde algún nodo inicial, rechazar.
                \State En caso contrario, aceptar.
                \end{algorithmic}
        \end{algorithm}

        Veamos que decide correctamente $L_b$. Si algún nodo marcado es alcanzable,
        entonces existe una entrada y una ejecución de $M$ que llega a una
        configuración en la que se ejecuta una transición con movimiento $I$. Por
        tanto, $M$ sí se mueve a la izquierda, y debemos rechazar
        $\langle M\rangle$. \\

        Recíprocamente, si ningún marcado es alcanzable, entonces ninguna transición
        con movimiento $I$ puede ejecutarse antes del primer movimiento a la
        izquierda. Pero todo movimiento a la izquierda tendría que ser precisamente
        el primero en algún instante. Por tanto, $M$ nunca se mueve a la izquierda
        en ninguna ejecución, y debemos aceptar $\langle M\rangle$.

        \item[c)] El conjunto de los pares $(M, w)$ tales que $M$, al actuar sobre la entrada $w$, nunca lee una casilla de la cinta más de una vez. \\
        
        Sea $$L_c = \{\langle M, w \rangle : M \text{ actuando sobre $w$ nunca lee una casilla de la cinta más de una vez}\}$$

        Como estamos hablando de ``una'' cinta, si el cabezal de dicha cinta cambia de sentido (iba hacia la derecha, 
        y vuelve a la izquierda, o iba hacia la izquierda y vuelve a la derecha), necesariamente vuelve a una casilla
        ya leída. \\
        
        Por tanto, mientras no se repita ninguna caislla, el movimiento tiene que ser siempre el mismo,
        o hacia la derecha, o hacia la izquierda (tampoco puede pararse). \\

        \begin{algorithm}
            \caption{MT $M_{L_c}$}
            \begin{algorithmic}
            \State \textbf{Entrada:} Una codificación $\langle M,w\rangle$
            \State Simulamos $M$ sobre $w$.
            \State Guardamos la posición actual del cabezal y el conjunto de casillas ya leídas.
            \State También guardamos la dirección del primer movimiento, si ya se ha producido.

            \While{la simulación no haya terminado}
                \If{la casilla actual ya había sido leída}
                    \State Rechazar
                \EndIf

                \State Marcamos la casilla actual como leída.

                \If{$M$ se para}
                    \State Aceptar
                \EndIf

                \State Ejecutamos un paso de $M$.

                \If{el movimiento cambia respecto al sentido anterior}
                    \State Rechazar
                \EndIf

                \If{la cabeza está fuera de la zona inicial de la palabra $w$, leyendo blancos nuevos, y se repite un estado moviéndose siempre en la misma dirección}
                    \State Aceptar
                \EndIf
            \EndWhile
            \end{algorithmic}
        \end{algorithm}

        El algoritmo siempre termina porque:
        \begin{enumerate}
            \item Si $M$ vuelve a leer una casilla, lo detectamos y rechazamos.
            \item Si $M$ se para antes de repetir casilla, aceptamos.
            \item Si $M$ sigue indefinidamente sin repetir casillas, 
            entonces necesariamente avanza siempre en una dirección leyendo blancos nuevos, 
            y como hay finitos estados, algún estado debe repetirse y entonces aceptamos.
        \end{enumerate}

    \end{itemize}
\end{ejercicio}



% Ejercicio 15

\begin{ejercicio}
    Indicar si los siguientes problemas son decidibles, semidecidibles o no semidecidibles:
    \begin{itemize}
        \item[a)] (MANUAL) Determinar si una MT se para para cualquier entrada. \\
        
        Hecho en el manual, sabemos que es no semidecidible. Se usa una reducción desde C-UNIVERSAL.

        \item[b)] Determinar si una MT no se para para ninguna entrada. \\

        El complementario de este problema es semidecidible con una MTND que seleccione no deterministícamente
        una entrada y simule la MT sobre la palabra. Si alguna para, acepta. Dicho complementario no es decidible 
        porque puede reducirse desde PARADA$(M,w)$ (simulamos $M$ sobre $w$. Si $M$ para, entonces
        la salida para también. Si $M$ no para, tampoco lo hace la salida por el algoritmo), que no es decidible. 
        Por tanto, el problema será no semidecidible. 

        \item[c)] Determinar si una MT se para, al menos, para una entrada. \\
        
        Discutido en el apartado anterior. Es el complementario mencionado (semidecidible pero no decidible).

        \item[d)] Determinar si una MT no se para, al menos, para una entrada. \\
        
        Sabemos que PARADA es semidecidible pero no decidible, luego su complementario es no semidecidible. Podemos
        reducir desde C-PARADA$(M,w)$ para ver que este problema no es semidecidible (simulamos $M$ sobre $w$. 
        Si para, la transformación algorítmica para, y si no para, tampoco para la salida del algoritmo).

    \end{itemize}
\end{ejercicio}



% Ejercicio 16

\begin{ejercicio}
    Supongamos que tenemos $4$ lenguajes $L_1$, $L_2$, $L_3$, $L_4$ de tal manera que existe una reducción de $L_1$ a $L_2$, de $L_2$ a $L_3$ y de $L_4$ a $L_3$. Para cada una de las siguientes afirmaciones,
    indicar si para todos los casos son CIERTAS, para todos los casos son FALSAS o si son POSIBLES (en algunos casos son ciertas y, en otros, son falsas).
    \begin{itemize}
        \item[a)] $L_1$ es recursivamente enumerable pero no recursivo, y $L_3$ es recursivo. \\
        
        \boxed{\textbf{FALSA}}. Si $L_3$ fuera recursivo, entonces también lo sería $L_2$, ya que $L_2 \propto L_3$, y por el 
        mismo argumento, también sería recursivo $L_1$, pues $L_1 \propto L_2$, y en particular $L_1$ sería r.e., lo cual es una contradicción.

        \item[b)] $L_1$ no es recursivo y $L_4$ no es recursivamente enumerable. \\
        
        \boxed{\textbf{POSIBLE}}. Consideramos $L_u$ (lenguaje asociado a UNIVERSAL), que es r.e. pero no recursivo, 
        y $L_d$ (lenguaje asociado a DIAGONAL), que no es r.e., luego tampoco recursivo. \\

        Ejemplo: $L_i = L_d \quad i = 1, \ldots, 4$. Todas las reducciones se dan por la identidad, y 
        $L_1 = L_d$ no es recursivo, y $L_4 = L_d$ no es r.e. \\

        Contraejemplo: $L_i = L_u \quad i = 1, \ldots, 4$. Todas las reducciones se dan por la identidad, pero ahora
        $L_1 = L_u$ no es recursivo, pero $L_4 = L_u$ sí es r.e. 

        \item[c)] El complementario de $L_1$ no es recursivamente enumerable, pero el complementario de $L_2$ es 
        recursivamente enumerable. \\

        \boxed{\textbf{FALSA}}. Sabemos que $L_1 \propto L_2 \iff \overline{L_1} \propto \overline{L_2}$. Entonces,
        si $\overline{L_2}$ fuera r.e., también lo sería $\overline{L_1}$. Contradicción.  

        \item[d)] El complementario de $L_2$ no es recursivo, pero el complementario de $L_3$ es recursivo. \\
        
        \boxed{\textbf{FALSA}}. Igual que en el apartado anterior, $L_2 \propto L_3 \iff \overline{L_2} \propto \overline{L_3}$. 
        Entonces, si $\overline{L_3}$ fuera recursivo, también lo sería $\overline{L_2}$. Contradicción. 

        \item[e)] Si $L_1$ es recursivo, entonces el complementario de $L_2$ es recursivo. \\
        
        \boxed{\textbf{POSIBLE}}. Consideramos $L_{V} = \emptyset$ (lenguaje vacío), que es recursivo, luego r.e., 
        y $A^*$ (todas las palabras del alfabeto $A = \{0,1\}$), que es recursivo. \\

        Ejemplo: $L_i = L_V \quad i = 1, \ldots, 4$. Todas las reducciones se dan por la identidad, y 
        $L_1 = L_V$ es recursivo, y $\overline{L_2} = A^{*}$ es recursivo. \\

        Contraejemplo: $L_1 = A^{*}$ y $L_i = L_u \quad i = 2, \ldots, 4$. Se reduce $A^{*} = L_1 \propto L_2 = L_u$ 
        enviando toda palabra de $L_1$ a una palabra fija de $L_2$. Las reducciones $L_2 \propto L_3$ y $L_4 \propto L_3$
        se dan por la identidad, pero ahora $L_1 = A^{*}$ es recursivo, pero $\overline{L_2} = \overline{L_u}$, que 
        no es recursivo (porque está asociado a C-UNIVERSAL, que no es semidecidible, luego $\overline{L_u}$ no es r.e., y tampoco
        recursivo). 

        \item[f)] Si $L_3$ es recursivo, entonces el complementario de $L_4$ es recursivo. \\
        
        \boxed{\textbf{CIERTA}}. Si $L_3$ es recursivo, como $L_4 \propto L_3$, $L_4$ es recursivo, y entonces, por
        un teorema visto en teoría, $\overline{L_4}$ también.

        \item[g)] Si $L_3$ es recursivamente enumerable, entonces la unión de $L_2$ y $L_4$ es recursivamente enumerable. \\
        
        \boxed{\textbf{CIERTA}}. Si $L_3$ es r.e., como $L_2 \propto L_3$ y $L_4 \propto L_3$, se tiene que 
        $L_2$ y $L_4$ son r.e., y como la unión de dos r.e. es r.e., entonces la afirmación es cierta.

        \item[h)] Si $L_3$ es recursivamente enumerable, entonces la intesección de $L_2$ y $L_4$ es recursivamente enumerable. \\
        
        \boxed{\textbf{CIERTA}}. Ya sabemos que si $L_3$ es r.e. entonces $L_2$ y $L_4$ son r.e., y como la 
        intersección de dos r.e. también es r.e., la afirmación es cierta.
    \end{itemize}
\end{ejercicio}



% Ejercicio 17

\begin{ejercicio}
    Sea el problema de determinar si una máquina de Turing acepta a lo más 100 palabras. 
    Determinar si es decidible, semidecidible o no semidecidible. Justificar la respuesta. \\

    Sea $P(M) = \{\langle M \rangle : |L(M)| \leq 100\}$. Veamos que es no semidecidible. 
    Para ello, consideramos el problema complementario $\overline{P}(M) = \{\langle M \rangle : |L(M)| > 100\}$. \\
    
    La propiedad ``$|L(M)| > 100$'' solo depende de $L(M)$, no de la MT $M$ concreta. Por tanto, es una propiedad
    de los lenguajes r.e.. Además, existe una MT que no la cumple (por ejemplo, una MT sin estados finales), y 
    existe una MT que la cumple (por ejemplo, una MT con mismo estado inicial y final). Por lo tanto, es no trivial,
    y entonces, por el Teorema de Rice, $\overline{P}$ es no decidible. \\
    
    Para la semidecidibilidad, usamos una MTND descrita como sigue:
    
    \begin{algorithm}
        \caption{MTND $M_{P}$}
        \label{algorithm:al-menos-101-semidecidible}
        \begin{algorithmic}
            \State \textbf{Entrada:} Una codificación $\langle M \rangle$ de una MT $M$

            \State Selecciona de forma no determinista $101$ palabras $u_1,\dots,u_{101} \in \{0,1\}^{*}$.

            \If{$\exists i,j \in \{1,\dots,101\} : i\neq j \quad \land \quad u_i=u_j$}
                \State $M_{P}$ rechaza.
            \EndIf

            \For{$i=1, \ldots, 101$}
                \State Simulamos $M$ con entrada $u_i$.
            \EndFor

            \If{todas las simulaciones aceptan}
                \State $M_{P}$ acepta.
            \EndIf
        \end{algorithmic}
    \end{algorithm}

    Tenemos entonces que, si $M$ acepta al menos $101$ palabras distintas 
    $u_1,\dots,u_{101}$, existe una rama de cómputo de $M_P$ que precisamente 
    toma $u_1,\dots,u_{101}$, comprueba que son distintas dos a dos, simula 
    $M(u_1),\dots,M(u_{101})$, y como todas aceptan, esa rama acepta. \\

    Por otro lado, si $M$ no acepta al menos $101$ palabras distintas, entonces 
    toda rama de cómputo de $M_P$ o bien selecciona dos palabras iguales, en cuyo 
    caso rechaza, o bien selecciona $101$ palabras distintas. En este último caso, 
    al menos una de ellas no pertenece a $L(M)$, luego al menos una simulación no 
    aceptará y, consecuentemente, esa rama no aceptará. Por lo tanto, $M_P$ acepta 
    exactamente las codificaciones $\langle M\rangle$ tales que $M$ acepta al menos 
    $101$ palabras distintas, de donde $\overline{P}$ es semidecidible. \\

    Como $\overline{P}$ es semidecidible pero no decidible, su complementario $\overline{\overline{P}} = P$ es 
    no semidecidible.
\end{ejercicio}



% Ejercicio 18

\begin{ejercicio}\label{ejercicio:ej18}
    Determinar cuáles de los siguientes problemas son decidibles, semidecidibles o no semidecidibles:
    \begin{itemize}
        \item[a)] Saber si una MT para una entrada 0011 no va a usar más de 10 casillas de la cinta.\\
        
        Veamos que es decidible, para ello, siguiendo el manual, basta ver que ``el problema se reduce a realizar un número
        finito de comprobaciones, cada una de ellas en un número finito de pasos.'' Definimos 
        $$\text{P}(M) = \{\langle M \rangle : \text{ con entrada 0011, no usa más de 10 casillas}\}$$
        Primero, la propiedad ``con entrada 0011 no usa más de 10 casillas'' no es una propiedad
        del lenguaje $L(M)$, luego no es una propiedad de los lenguajes r.e., y no puede aplicarse el Teorema de Rice. Para ver que es decidible, 
        consideramos una MT $M_{P}$ que funciona como sigue:
        \begin{algorithm}
            \caption{MT $M_{P}$}
            \label{algorithm:parada0011-decidible}
            \begin{algorithmic}
                \State \textbf{Entrada:} Una codificación $\langle M \rangle$ de una MT $M$

                \State Simulamos $M$ con entrada $0011$, partiendo de su configuración inicial.
                \State En una segunda cinta vamos guardando las casillas distintas visitadas por $M$ y su número (contador).

                \State \hspace{1cm} \textbf{Si} contador $ > 10$ \textbf{entonces}

                \State \hspace{2cm} $M_P$ rechaza.

                \State \hspace{1cm} \textbf{Si} $M$ se para y contador $\leq 10$ \textbf{entonces}

                \State \hspace{2cm} $M_P$ acepta.

                \State \hspace{1cm} \textbf{Si} se repite una configuración y contador $\leq 10$ \textbf{entonces}

                \State \hspace{2cm} $M_P$ acepta.
            \end{algorithmic}
        \end{algorithm}

        Al simular $M_p$, pueden pasar tres cosas:
        \begin{enumerate}
            \item Que intente usar más de 10 casillas, y entonces $M_P$ debe rechazar.
            \item Que pare sin superar 10 casillas, y entonces $M_P$ debe aceptar.
            \item Que no pare (cicle), pero entonces, como hay finitas
            configuraciones posibles, necesariamente se repetirá alguna en tiempo 
            finito, y entonces ciclará siempre sin salir de esas 10 casillas, y entonces
            $M_P$ debe aceptar.
        \end{enumerate}

        Por lo tanto, $M_P$ es un algoritmo que resuelve el problema, lo cual lo hace decidible.

        

        \item[b)] Dadas dos MT saber si aceptan el mismo lenguaje. \\ 
        
        Definimos $\text{IGUAL}(M_1, M_2) = \{(\langle M_1 \rangle, \langle M_2 \rangle) : L(M_1) = L(M_2)\}$. Veamos 
        que no es semidecidible. Podría hacerse directamente, pero primero, para ver que no es decidible, 
        fijamos una MT $M_{\emptyset}$ tal que $$L(M_{\emptyset})=\emptyset.$$ Consideramos el problema
        $$P_{M_{\emptyset}}=\{\langle M\rangle : L(M)=L(M_{\emptyset})\}.$$ Como $L(M_{\emptyset})=\emptyset$, 
        entonces $$L(M)=L(M_{\emptyset}) \iff L(M)=\emptyset.$$ Por tanto, 
        $$P_{M_{\emptyset}}=\{\langle M\rangle : L(M)=\emptyset\}.$$
        Si IGUAL fuera decidible, también lo sería $P_{M_{\emptyset}}$, lo cual contradice el Teorema de Rice, 
        ya que la propiedad ``$L(M)=\emptyset$'' es una propiedad de los lenguajes r.e. (depende solo de $L(M)$) y 
        además es no trivial: la cumple una MT sin estados finales, pero no la cumple una MT que acepte, por ejemplo, 
        la palabra $0$. Por tanto, IGUAL no es decidible. \\

        Veamos que, además, no es semidecidible por medio de una reducción desde
        $$\text{VACIO}(M) = \{\langle M \rangle : L(M) = \emptyset\}$$

        Fijada una MT $M$, tal que $L(M) = \emptyset$, es decir, $M \in \text{VACIO}$, definimos la reducción 
        $$F(\langle N \rangle) = (\langle N \rangle, \langle M \rangle)$$
        Intuitivamente, comparar una MT fijada vacía convierte VACIO en IGUAL.
        Entonces:
        \begin{enumerate}
            \item Si $L(N) = \emptyset$, se tiene que $L(N) = \emptyset = L(M)$, luego $F(\langle N \rangle) \in \text{IGUAL}$.
            \item Si $L(N) \neq \emptyset$, se tiene que $L(N) \neq L(M)$, luego $F(\langle N \rangle) \notin \text{IGUAL}$.
        \end{enumerate}
        Se tiene entonces que $\text{VACIO} \propto \text{IGUAL}$. Como $\text{VACIO}$ no es semidecidible, tampoco lo es $\text{IGUAL}$.
    \end{itemize}
\end{ejercicio}



% Ejercicio 19

\begin{ejercicio}
    Determinar cuáles de los siguientes problemas son decidibles, semidecidibles o no semidecidibles:
    \begin{itemize}
        \item[a)] Dada una MT $M$ y una palabra $u$, saber si la MT acepta la palabra $u$ en un número de pasos menor o igual a $|u|$. \\
        
        Definimos $$P(M,u)=\{\langle M,u\rangle : M \text{ acepta } u \text{ en un número de pasos menor o igual a } |u|\}.$$

        Veamos que es decidible. Para ello, consideramos una MT $M_P$ que funciona como sigue, donde la configuración inicial de una MT con entrada $u$ es $(q_0,\varepsilon,u)$:

        \begin{algorithm}
        \caption{MT $M_{P}$}
        \label{algorithm:acepta-u-en-leq-longitud}
        \begin{algorithmic}
            \State \textbf{Entrada:} Una codificación $\langle M,u \rangle$ de una MT $M$ y una palabra $u$

            \State \textbf{Para} $i = 0, \ldots, |u|$ \textbf{hacer}

                \State \hspace{1cm} Simulamos $i$ pasos de $M$ con entrada $u$, partiendo de su configuración inicial.

                \State \hspace{1cm} \textbf{Si} $M$ entra en un estado de aceptación \textbf{entonces} 
                
                \State \hspace{2cm} $M_P$ acepta.

                \State \hspace{1cm} \textbf{Si} $M$ para sin aceptar \textbf{entonces} 

                \State \hspace{2cm} $M_P$ rechaza.

                \State \hspace{1cm} \textbf{Si} $i = |u|$ y $M$ no acepta \textbf{entonces} 
                
                \State \hspace{2cm} $M_P$ rechaza.
        \end{algorithmic}
        \end{algorithm}

        Al simular $M_P$, pueden pasar tres cosas:

        \begin{enumerate}
            \item Que $M$ acepte en $k$ pasos con $k \le |u|$, y entonces $M_P$ debe aceptar.

            \item Que $M$ se pare antes o en $|u|$ pasos, pero sin aceptar, y entonces $M_P$ debe rechazar.

            \item Que $M$ no haya aceptado tras $|u|$ pasos. En ese caso, aunque después aceptara o ciclara, ya no cumpliría la condición pedida, y entonces $M_P$ debe rechazar.
        \end{enumerate}

        Por tanto, $M_P$ siempre termina, ya que sólo necesita simular, como mucho, $|u|$ pasos de cálculo, y decide correctamente si $M$ acepta $u$ en un número de pasos menor o igual a $|u|$,
        y el problema es decidible.

        

        \item[b)] Dada una MT determinar si no acepta la palabra vacía.\\
        
        Deberíamos intuir, al ser de la forma ``determinar si no acepta (algo)'' que no es semidecidible. Definimos $P_\varepsilon = \{\langle M \rangle : M \text{ no acepta } \varepsilon\}$. Vemos que es el complementario
        de $C-P_\varepsilon = \{\langle M \rangle : M \text{ acepta } \varepsilon\}$. Sin embargo, $C-P_\varepsilon$ es 
        semidecidible, ya que podemos construir una MT $M_{P_\varepsilon}$ que, con entrada $\langle M \rangle$, 
        simule a $M$ sobre la palabra vacía $\varepsilon$, y acepte si $M$ acepta. Si además $P_\varepsilon$ fuera
        semidecidible entonces $P_\varepsilon$ y $C-P_\varepsilon$ serían semidecidibles, lo que haría que 
        $P_\varepsilon$ fuera decidible, y a su vez, que $C-P_\varepsilon$, al ser complementario de decidible, también lo fuera.
        Sin embargo, $C-P_\varepsilon$ no es decidible, ya que ``aceptar $\varepsilon$'' depende solo de $L(M)$, porque
        equivale a ver si $\varepsilon \in L(M)$. Entonces es una propiedad de los lenguajes r.e., y además es no trivial, 
        ya que hay una MT cuyo lenguaje aceptado es $\{\varepsilon\}$ (por ejemplo, una transición que lea blanco y pase a estado final), que sí la cumple, 
        pero también otra MT cuyo lenguaje aceptado es $\{0\}$, que no la cumple. Por tanto, $C-P_\varepsilon$
        no es decidible por el Teorema de Rice, lo cual sería una contradicción. Necesariamente entonces $P_\varepsilon$ no es semidecidible.
    \end{itemize}
\end{ejercicio}



% Ejercicio 20

\begin{ejercicio}
    Determinar cuáles de los siguientes problemas son decidibles, semidecidibles o no semidecidibles 
    (se supone que las MTs tienen a $A = \{0, 1\}$ como alfabeto de entrada):
    \begin{itemize}
        \item[a)] Dadas dos máquinas de Turing, $M_1$ y $M_2$, determinar si existe, al menos, una palabra
        aceptada simultáneamente por ambas máquinas. \\

        Sea $P = \{\langle M_1, M_2 \rangle : L(M_1) \cap L(M_2) \neq \emptyset\}$. Veamos que es semidecidible pero no 
        decidible. La no decidibilidad se deduce del Teorema de Rice, ya que, fijada una MT $M_{*}$ que acepte 
        todas las palabras, $L(M_*) = A^{*}$, consideramos el problema siguiente 
        $$P_{M_*} = \{\langle M \rangle : L(M) \cap L(M_{*}) \neq \emptyset\}$$ Como $L(M_*) = A^{*}$, entonces
        $$L(M) \cap L(M_{*}) \neq \emptyset \iff L(M) \neq \emptyset$$ Por tanto,
        $P_{M_*} = \{\langle M \rangle : L(M) \neq \emptyset\}$. Si $P$ fuera decidible, también lo sería $P_{M_*}$,
        lo cual es una contradicción con el Teorema de Rice, ya que la propiedad ``$L(M) \neq \emptyset$'' es una propiedad 
        de los lenguajes r.e. (solo depende de $L(M)$) y además es no trivial, ya que si $M$ es una MT sin estados finales, 
        entonces $L(M) = \emptyset$, y si $M$ es una MT que acepta todas las palabras, se tiene que 
        $L(M) = A^{*} \neq \emptyset$. Por tanto, $P$ es no decidible. \\

        La semidecidibilidad se obtiene con la siguiente MTND:

        \begin{algorithm}
            \caption{MTND $M_{P}$}
            \label{algorithm:palabra-simultanea-semidecidible}
            \begin{algorithmic}
                \State \textbf{Entrada:} Un par de codificaciones $\langle M_1, M_2 \rangle$ de dos MT $M_1$ y $M_2$

                \State Selecciona de forma no determinista una palabra $w \in A^{*}$.

                \For{$i=1,2$}
                    \State Simulamos $M_i$ con entrada $w$.
                \EndFor

                \If{ambas simulaciones aceptan}
                    \State $M_{P}$ acepta.
                \EndIf
            \end{algorithmic}
        \end{algorithm}

        \item[b)] (MANUAL) Dada una MT, determinar si para toda palabra de entrada no realiza más de 5 movimientos. \\
        
        Hecho en el manual con 10 movimientos (pasos de cálculo). La clave está en ver que en 5 movimientos solo se pueden
        leer cinco casillas de la cinta como máximo. Entonces las palabras se identifican para este problema con los 
        prefijos de longitud, a lo sumo, 5. Así, dada una palabra $u \in A^{*}$ tal que $|u| \leq 5$, y dada 
        una palabra $v \in A^{*}$, entonces, para una MT $M$ cualquiera, 
        $$M \text{ no realiza más de 5 movimientos con $u$} \iff M \text{ no realiza más de 5 movimientos con $uv$}$$

        El conjunto $\{u \in A^{*} : |u| \leq 5\}$ es finito. De hecho, podemos obtener su cardinal:

        $$\sum_{k=0}^{5} 2^k = \dfrac{2^{6} - 1}{2 - 1} = 63$$

        y basta ejecutar 5 pasos de cálculo para cada una de estas 63 palabras. Si con todas para, la respuesta es que sí, 
        y si con alguna no ha parado, la respuesta es que no.

        \item[c)] Determinar si una MT no es determinista. \\
        
        Sea $P = \{\langle M \rangle : M \text{ es una MTND}\}$. Podemos utilizar el siguiente algoritmo para resolver el problema:

        \begin{algorithm}
            \caption{MT $M_{P}$}
            \label{algorithm:no-determinista-decidible}
            \begin{algorithmic}
                \State \textbf{Entrada:} Una codificación $\langle M \rangle$ de una MT $M$

                \State Leemos la codificación de $M$, en particular su conjunto de estados $Q$, su alfabeto de cinta $B$ y su tabla de transiciones.

                \For{cada par $(q,a) \in Q \times B$}
                    \State Contamos cuántas transiciones están definidas desde $(q,a)$.
                    
                    \If{hay dos o más transiciones definidas desde $(q,a)$}
                        \State $M_{P}$ acepta.
                    \EndIf
                \EndFor

                \State $M_{P}$ rechaza.
            \end{algorithmic}
        \end{algorithm}

        Como $Q$, $B$ y la tabla de transiciones de $M$ son finitos, el algoritmo siempre termina, y el algoritmo 
        anterior responde correctamente en todos los casos. Por tanto, $P$ es decidible y, en particular, semidecidible.

        \item[d)] (MANUAL) Dada una MT, determinar si ninguna de las palabras que acepta es un palíndromo. \\
        
        Es el complementario del problema del Ejercicio \ref{ejercicio:ej9}, que es semidecidible pero no decidible, luego este problema
        es no semidecidible.

    \end{itemize}
\end{ejercicio}



% Ejercicio 21

\begin{ejercicio}
    Se considera el siguiente problema: dadas dos máquinas de Turing, determinar si aceptan el mismo lenguaje. 
    Razonar si este problema es decidible, semidecidible o no semidecidible. \\

    Es igual que el Ejercicio \ref{ejercicio:ej18} b).
\end{ejercicio}



% Ejercicio 22

\begin{ejercicio}
    Determinar, justificando las respuestas, cuáles de los siguientes problemas son decidibles, semidecidibles o no semidecidibles:
    \begin{itemize}
        \item[a)] Dada una máquina de Turing y una palabra de entrada determinar si visita menos de 10 casillas 
        distintas de la cinta de lectura. \\ % Decidible 

        Es igual que el Ejercicio \ref{ejercicio:ej18} a), pero con una entrada $w \in \{0,1\}^{*}$ cualquiera. \\

        Solo cambiaría la definición del problema $$P(M,w) = \{\langle M, w \rangle : M \text{ con entrada $w$ visita menos de 10 casillas distintas}\}$$

        y la entrada del algoritmo sería una codificación $\langle M,w \rangle$ de una MT $M$ y una palabra $w$, y habría
        que simular $M$ sobre $w$ y seguir la misma lógica que en dicho ejercicio. \\

        Este problema también es decidible.

        \item[b)] Dadas dos Máquinas de Turing, determinar si el conjunto de las palabras aceptadas a la vez por 
        ambas máquinas de Turing es finito. \\ % No semidecidible

        Podríamos ver en dos pasos no decidibilidad con Rice, y luego no semidecidibilidad, pero veamos directamente
        que no es semidecidible por medio de una reducción desde VACIO (que sabemos por teoría que no es semidecidible). \\

        Definimos $P(M_1, M_2) = \{\langle M_1, M_2 \rangle : L(M_1) \cap L(M_2) \text{ es finito}\}$. Reducimos desde
        $$\text{VACIO}(M) = \{\langle M \rangle : L(M) = \emptyset\}$$

        Dada una MT $M$, construimos $M'$ dada como sigue:
        
        \begin{algorithm}
            \caption{MT $M'$}
            \label{algorithm:N-prima}
            \begin{algorithmic}
                \State \textbf{Entrada:} Una palabra $x \in \{0,1\}^{*}$

                \State Ignoramos la entrada $x$.

                \For{$t = 1,\dots, \infty$}
                    \For{cada palabra $w \in \{0,1\}^{*}$ con $|w| \leq t$}
                        \State Simulamos $M$ con entrada $w$ durante $t$ pasos.
                        \If{$M$ acepta $w$ en esos $t$ pasos}
                            \State $M'$ acepta $x$.
                        \EndIf
                    \EndFor
                \EndFor
            \end{algorithmic}
        \end{algorithm}

        Es decir, se tiene que $$L(M') = \begin{cases}
            \emptyset &\text{si } L(M) = \emptyset \\
            \{0,1\}^{*} &\text{si } L(M) \neq \emptyset 
        \end{cases}$$

        Fijamos ahora una MT $M_0$ tal que $L(M_0) = \{0,1\}^{*}$. Definimos la reducción
        $$F(\langle M \rangle) = (\langle M' \rangle, \langle M_0 \rangle)$$

        Entonces $L(M') \cap L(M_0) = L(M')$, y 
        $$L(M) = \emptyset \iff L(M') = \emptyset \iff L(M') \cap L(M_0) \text{ es finito}$$

        Por tanto, $\text{VACIO} \propto P$. Como VACIO no es semidecidible, tampoco lo es $P$. 


    \end{itemize}
\end{ejercicio}



% Ejercicio 23

\begin{ejercicio}
    Indica cuáles de los siguientes problemas son decidibles, semidecidibles o no semidecidibles:
    \begin{itemize}
        \item[a)] Determinar si el lenguaje aceptado por una MT es regular. \\ % No decidible
        
        Sea $\text{REGULAR} = \{ \langle M \rangle : L(M) \text{ es regular}\}$. Podemos ver que no es decidible y luego que 
        no es semidecidible en dos pasos, o bien directamente que no es semidecidible por medio de una reducción 
        desde VACIO. Para ver que es no decidible usamos el Teorema de Rice, ya que la propiedad ``$L(M)$ es regular'' 
        depende solo de $L(M)$, luego es una propiedad de los lenguajes r.e. y es no trivial, porque 
        existe una MT $M_1$ con $L(M_1) = \emptyset$ (sin estados finales por ejemplo), y $L(M_1)$  cumple la propiedad, 
        y existe una MT $M_2$ con $L(M_2) = \{0^n 1^n : n \geq 0\}$ (por ejemplo, multicinta que cuente cuántos 0's
        aparecen y ponga ese mismo número de 1's a continuación), y $L(M_2)$ no cumple la propiedad. \\

        La no semidecidibilidad puede obtenerse con una reducción desde 
        $$\text{VACIO}(M) = \{\langle M \rangle : L(M) = \emptyset\}$$

        Dada una MT $M$, construimos una MT M' con el siguiente algoritmo:

        \begin{algorithm}
            \caption{MT $M'$}
            \label{algorithm:regular-prima}
            \begin{algorithmic}
                \State \textbf{Entrada:} Una palabra $x \in \{0,1\}^{*}$

                \State Comprobamos si $x \in \{0^n1^n : n \geq 0\}$.

                \If{$x \notin \{0^n1^n : n \geq 0\}$}
                    \State $M'$ rechaza.
                \EndIf

                \For{$t = 1,\dots,\infty$}
                    \For{cada palabra $w \in \{0,1\}^{*}$ con $|w| \leq t$}
                        \State Simulamos $M$ con entrada $w$ durante $t$ pasos.
                        \If{$M$ acepta $w$ en esos $t$ pasos}
                            \State $M'$ acepta $x$.
                        \EndIf
                    \EndFor
                \EndFor
            \end{algorithmic}
        \end{algorithm}
        
        Tendremos entonces que $$L(M') = \begin{cases}
            \emptyset &\text{si } L(M) = \emptyset \\
            \{0^n 1^n : n \geq 0\} &\text{si } L(M) \neq \emptyset 
        \end{cases}$$

        Como $\emptyset$ es regular, pero $\{0^n 1^n : n \geq 0\}$ no lo es, se tiene que 

        $$L(M) = \emptyset \iff L(M') = \emptyset \iff L(M') \text{ es regular}$$

        Definimos la reducción $F(\langle M \rangle) = \langle M' \rangle$. Viendo que

        $$\langle M \rangle \in \text{VACIO} \iff F(\langle M \rangle) = \langle M' \rangle \in \text{REGULAR}$$

        deducimos que $\text{VACIO} \propto \text{REGULAR}$. Como \text{VACIO} no es semidecidible, tampoco lo es 
        \text{REGULAR}.

        \item[b)] Dadas dos MTs determinar si hay una palabra que es aceptada por las dos MT y 
        además es un palíndromo. \\ % Semidecidible y no decidible

        Sea $P = \{\langle M_1, M_2 \rangle : \exists w \in \{0,1\}^{*} \text{ palíndromo} : w \in L(M_1) \cap L(M_2)\}$. 
        Veamos que es semidecidible pero no decidible. La no decidibilidad se deduce del Teorema de Rice, ya que, 
        fijada una MT $M_{*}$ que acepte todas las palabras, $L(M_*) = A^{*}$, consideramos el problema siguiente 
        $$P_{M_*} = \{\langle M \rangle : \exists w \in \{0,1\}^{*} \text{ palíndromo} : w \in L(M) \cap L(M_*)\}$$ 
        Como $L(M_*) = A^{*}$, entonces, considerando $\text{PAL} = \{w \in \{0,1\}^{*} : w \text{ es un palíndromo}\}$:
        $$\exists w \in \{0,1\}^{*} \text{ palíndromo} : w \in L(M) \cap L(M_*) \iff L(M) \cap \text{PAL} \neq \emptyset$$ 
        Por tanto, $P_{M_*} = \{\langle M \rangle : L(M) \cap \text{PAL} \neq \emptyset\}$. Si $P$ fuera decidible, 
        también lo sería $P_{M_*}$, lo cual es una contradicción con el Teorema de Rice, ya que la propiedad 
        ``$L(M) \cap \text{PAL} \neq \emptyset$'' es una propiedad de los lenguajes r.e. (solo depende de $L(M)$) y 
        además es no trivial, ya que si $M$ es una MT sin estados finales, entonces $L(M) = \emptyset$, y no cumple
        la propiedad, y si $M$ es una MT que acepta todas las palabras, se tiene que 
        $L(M) = A^{*} \neq \emptyset$, que sí cumple la propiedad (contiene a todos los palídromos de hecho). 
        Por tanto, $P$ es no decidible. \\

        La semidecidibilidad se obtiene con la siguiente MTND:

        \begin{algorithm}
            \caption{MTND $M_{P}$}
            \label{algorithm:palindromo-simultaneo-semidecidible}
            \begin{algorithmic}
                \State \textbf{Entrada:} Un par de codificaciones $\langle M_1, M_2 \rangle$ de dos MT $M_1$ y $M_2$

                \State Selecciona de forma no determinista una palabra $w \in A^{*}$.

                \State Comprueba si $w$ es un palíndromo.

                \If{$w$ no es un palíndromo}
                    \State $M_p$ rechaza.
                \EndIf

                \For{$i=1,2$}
                    \State Simulamos $M_i$ con entrada $w$.
                \EndFor

                \If{ambas simulaciones aceptan}
                    \State $M_{P}$ acepta.
                \EndIf
            \end{algorithmic}
        \end{algorithm}

    \end{itemize}
\end{ejercicio}



% Ejercicio 24

\begin{ejercicio}
    Indica cuáles de los siguientes problemas son decidibles, semidecidibles o no semidecidibles:
    \begin{itemize}
        \item[a)] Dada una MT determinar si todos sus estados son accesibles 
        (se puede llegar a ellos para un cálculo para alguna entrada). \\

        Sea $$P = \{\langle M \rangle : \forall q \in Q_M, \exists w \in \{0,1\}^{*} : \exists t \in \mathbb{N} : M(w) 
        \text{ accede al estado } q \text{ en } t \text{ pasos} \}$$

        El problema es semidecidible pero no decidible. Podemos ver la semidecidibilidad con la siguiente MTND:

        \begin{algorithm}
            \caption{MTND $M_{P}$}
            \label{algorithm:estados-accesibles-semidecidible}
            \begin{algorithmic}
                \State \textbf{Entrada:} Una codificación $\langle M\rangle$ de una MT $M$

                \State Sea $Q_M=\{q_1,\dots,q_k\}$ el conjunto de estados de $M$.

                \For{$i=1,\dots,k$}
                    \State Selecciona de forma no determinista una palabra $w_i\in\{0,1\}^{*}$.
                    \State Selecciona de forma no determinista un número $t_i\in\mathbb{N}$.
                \EndFor

                \For{$i=1,\dots,k$}
                    \State Simula $M$ con entrada $w_i$ durante, a lo sumo, $t_i$ pasos.
                    \If{durante esa simulación no se visita el estado $q_i$}
                        \State $M_p$ rechaza.
                    \EndIf
                \EndFor

                \State $M_p$ acepta.
            \end{algorithmic}
        \end{algorithm}

        Si todos los estados son accesibles, existirá, para cada $i = 1, \ldots, k$, un estado $q_i \in Q_M$, una palabra 
        $w_i \in \{0,1\}^{*}$ y un número $t_i$ tales que la ejecución de $M$ sobre $w_i$ visita $q_i$ en, a lo sumo,
        $t_i$ pasos. \\

        Para probar que $P$ no es decidible, reducimos desde
        $$\text{UNIVERSAL}(M,w)=\{\langle M,w\rangle : M \text{ acepta } w\},$$
        que sabemos que no es decidible.

        

        Dada una instancia $\langle M,w\rangle$ de UNIVERSAL, construimos una MT $M'$ tal que 
        todos sus estados sean accesibles salvo, quizá, un estado
        especial $q_*$. La máquina $M'$ ignora su entrada y funciona como sigue:

        \begin{algorithm}
            \caption{MT $M'$}
            \label{algorithm:accesibilidad-prima}
            \begin{algorithmic}
                \State \textbf{Entrada:} Una palabra $x\in\{0,1\}^{*}$

                \State $M'$ ignora la entrada $x$.

                \State Obtenemos $Q' \setminus \{q_*\}=\{p_0,p_1,\dots,p_m\}$, donde $p_0$ es el estado inicial de $M'$.

                \State $M'$ realiza una subrutina de accesibilidad que recorre todos los estados de $M'$ salvo $q_*$, 
                mediante transiciones auxiliares que no interfieren con la simulación posterior.

                \State Al llegar al estado $p_m$, $M'$ coloca en la cinta $\langle M \rangle 111 w$.

                \State Desde $p_m$, $M'$ comienza la simulación de $M$ con entrada $w$.

                \If{$M$ acepta $w$}
                    \State $M'$ entra en el estado especial $q_*$ y acepta.
                \EndIf

                \If{$M$ rechaza $w$}
                    \State $M'$ rechaza sin pasar por $q_*$.
                \EndIf

                \State Si $M$ no para con $w$, entonces $M'$ tampoco para y nunca entra en $q_*$.
            \end{algorithmic}
        \end{algorithm}

        Por construcción, todos los estados de $M'$ distintos de $q_*$ son accesibles, ya que
        la subrutina de accesibilidad pasa por todos ellos. \\

        El único estado cuya accesibilidad no queda garantizada entonces es
        $q_*$. Además, por construcción, $M'$ entra en $q_*$ si y solo si la simulación de
        $M$ sobre $w$ acepta. Deducimos de lo anterior que
        $$q_* \text{ es accesible en } M' \iff M \text{ acepta } w.$$

        Como todos los demás estados son accesibles por construcción, se tiene
        $$\text{todos los estados de }M'\text{ son accesibles} \iff q_* \text{ es accesible en } M' \iff
        M \text{ acepta } w.$$

        Por tanto, $$\langle M,w\rangle\in \text{UNIVERSAL} \iff \langle M'\rangle\in P. $$

        Así, $$\text{UNIVERSAL} \propto P.$$

        Si $P$ fuera decidible, entonces UNIVERSAL también sería decidible, contradicción.
        Por tanto, $P$ no es decidible.

        

        \item[b)] Dada una MT determinar si su número de estados es menor o igual a 5. \\
        
        Sea $P  = \{\langle M \rangle : |Q_M| \leq 5\}$. El problema es decidible, y se resuelve con el siguiente algoritmo:

        \begin{algorithm}
            \caption{MT $M_P$}
            \begin{algorithmic}
                \State \textbf{Entrada:} Una codificación $\langle M\rangle$ de una MT $M$
                \State Comprobamos que $\langle M\rangle$ es una codificación correcta de una MT.
                \State Leemos el conjunto de estados $Q$ de $M$.
                \State Contamos cuántos estados tiene $M$.
                \If{$|Q|\leq 5$}
                    \State $M_p$ acepta.
                \Else
                    \State $M_p$ rechaza.
                \EndIf
            \end{algorithmic}
        \end{algorithm}

    \end{itemize}
\end{ejercicio}



% Ejercicio 25

\begin{ejercicio}
    Demuestra que el Problema de la Correspondencia de Post con un alfabeto A que tiene un sólo elemento es decidible. \\

    Sea $a \in A$ tal que $A = \{a\}$. Entonces, $A^{*} = \{a^n : n \geq 0\}$. Cada bloque del PCP será de la forma 
    $$\dfrac{a^{p_i}}{a^{q_i}}$$ 
    con $p_i$ la longitud de la palabra superior, y $q_i$ la longitud de la palabra inferior, $i = 1, \ldots, k$, y 
    $k$ el número de bloques. Entonces, para una sucesión no vacía de enteros $i_1, \ldots, i_m$, la palabra de arriba 
    será $a^{p_{i_1}} \cdots a^{p_{i_m}} = a^{p_{i_1} + \cdots + p_{i_m}}$. Análogamente, la palabra de abajo será 
    $a^{q_{i_1} + \cdots + q_{i_m}}$. Entonces, como solo hay un símbolo, dos palabras son iguales si y solo si tienen
    la misma longitud, es decir, la sucesión es válida (como solución al PCP) si y solo si 
    $p_{i_1} + \cdots + p_{i_m} = q_{i_1} + \cdots + q_{i_m}$. Definimos $d_i = p_i - q_i$. La condición de solución
    se reescribe como
    $$p_{i_1} + \cdots + p_{i_m} = q_{i_1} + \cdots + q_{i_m} \iff 
    (p_{i_1} - q_{i_1}) + \cdots + (p_{i_m} - q_{i_m}) = 0 \iff$$$$ d_{i_1} + \cdots d_{i_m} = 0$$
    Para ver que es 
    decidible usamos el siguiente algoritmo:

    \begin{algorithm}
        \caption{MT $M_{PCP}$}
        \begin{algorithmic}
            \State \textbf{Entrada:} Una instancia de PCP sobre $A=\{a\}$ con bloques
            $\frac{a^{p_1}}{a^{q_1}},\dots,\frac{a^{p_k}}{a^{q_k}}$.

            \For{$i=1,\dots,k$}
                \State Calculamos $d_i = p_i-q_i$.
            \EndFor

            \If{existe $i$ tal que $d_i=0$}
                \State Aceptamos.
            \EndIf

            \If{existen $i,j$ tales que $d_i>0$ y $d_j<0$}
                \State Aceptamos.
            \EndIf

            \State Rechazamos.
        \end{algorithmic}
    \end{algorithm}

    Si algún $d_i = 0$, entonces el bloque $i$ es solución. Si hay $d_i > 0$ y $d_j < 0$ para $i,j \in \{1, \ldots, k\}$,
    escribimos $d_i = r > 0$, y $d_j = -s < 0$. Tomando $s$ copias del bloque $i$ y $r$ copias del bloque $j$,
    la diferencia total será $s \cdot r + r \cdot (-s) = 0$, luego hay solución. Si todos los $d_i$ son positivos,
    cualquier suma no vacía será positiva. Si, por otro lado, todos los $d_i$ son negativos, cualquier suma no vacía 
    será negativa. En ninguno de estos dos casos puede haber solución. Como el algoritmo solo recorre un número
    finito de bloques y siempre termina, el PCP sobre un alfabeto con un único símbolo es decidible.

\end{ejercicio}
